\subsection*{§. 196.\\ 椭圆几何中的滑移。}

现在专门讨论三维椭圆空间，研究该空间的运动群，正如 §. 189 已对双曲空间所作的那样。

采用由关系
\[
\tag{40}
 x_0^2+x_1^2+x_2^2+x_3^2=1,
\]
联系的 Weierstrass 坐标 $x_0,x_1,x_2,x_3$，则 $ds^2$ 为（§. 193）
\[
\tag{40*}
 ds^2=R^2(dx_0^2+dx_1^2+dx_2^2+dx_3^2).
\]
椭圆空间的每个运动，都由表示欧几里得空间的一个射影变换给出；在这个射影变换下，绝对二次曲面变为自身。因此，该运动由线性代换
\[
\tag{41}
\left\{
\begin{aligned}
 x'_0&=a_{00}x_0+a_{01}x_1+a_{02}x_2+a_{03}x_3,\\
 x'_1&=a_{10}x_0+a_{11}x_1+a_{12}x_2+a_{13}x_3,\\
 x'_2&=a_{20}x_0+a_{21}x_1+a_{22}x_2+a_{23}x_3,\\
 x'_3&=a_{30}x_0+a_{31}x_1+a_{32}x_2+a_{33}x_3.
\end{aligned}\right.
\]
表示。由于必须有
\[
 \mathop{\mathrm{S}}x_i'^2=\mathop{\mathrm{S}}x_i^2=1,
\]
显然这个代换必须是\emph{正交的}。所以，椭圆空间的连续运动群由全部正交代换 (41) 表示，其行列式必须为 $=+1$；亦即由四维欧几里得空间绕一个固定点（原点）的旋转表示。

这些椭圆空间运动中，有一类格外重要；双曲空间运动没有与之对应的类型，而它在许多方面类似于普通空间的平移群。把这类运动定义为：空间中所有点到运动后各自所占位置的距离都相等；称为\emph{滑移}（Schiebungen）。要证明它们存在并同时求出其解析表达式，只须根据第434页公式 (25) 表示出：由 (41) 算出的表达式
\[
 x_0x'_0+x_1x'_1+x_2x'_2+x_3x'_3,
\]
为常数。
立即得到
\[
 a_{00}=a_{11}=a_{22}=a_{33}=A=\cos\left(\frac\delta R\right),
\]
其中 $\delta$ 是空间中每一点移动的常长度，并且
\[
 a_{ik}+a_{ki}=0\quad (i\ne k).
\]
另一方面，(41) 必须是正交代换，由此得到
\[
 a_{00}a_{21}+a_{20}a_{11}=0,
 \qquad a_{30}^2+a_{00}^2=a_{21}^2+a_{31}^2,
\]
从而
\[
 a_{21}=\pm a_{30},\qquad a_{31}=\mp a_{20}.
\]
区分上号与下号两种情形，得到两类不同的滑移，分别由公式
\[
\tag{III}
\left\{
\begin{aligned}
 x'_0&=Ax_0-Bx_1-Cx_2-Dx_3,\\
 x'_1&=Bx_0+Ax_1-Dx_2+Cx_3,\\
 x'_2&=Cx_0+Dx_1+Ax_2-Bx_3,\\
 x'_3&=Dx_0-Cx_1+Bx_2+Ax_3,
\end{aligned}\right.
\]
及另一组公式
\[
\tag{III*}
\left\{
\begin{aligned}
 x'_0&=Ax_0-Bx_1-Cx_2-Dx_3,\\
 x'_1&=Bx_0+Ax_1-Dx_2-Cx_3,\\
 x'_2&=Cx_0-Dx_1+Ax_2+Bx_3,\\
 x'_3&=Dx_0+Cx_1-Bx_2+Ax_3,
\end{aligned}\right.
\]
给出；两处的 $A,B,C,D$ 都是满足关系
\[
 A^2+B^2+C^2+D^2=1.
\]
的任意实常数。显然，这四个常数可以解释为：\emph{点 $(1,0,0,0)$ 在滑移后所到达点的坐标。}

可见，当固定点 $(1,0,0,0)$ 或空间中任意另一个固定点在滑移后所应占的位置时，第一类滑移 (III) 或第二类滑移 (III*) 便完全确定。

考察任意一个滑移，例如第一类滑移 (III)。在此滑移下，每个平面 $(\xi_0,\xi_1,\xi_2,\xi_3)$ 变为新平面 $(\xi'_0,\xi'_1,\xi'_2,\xi'_3)$，而平面坐标的变换公式与 (III) 相同，即
\[
\begin{aligned}
 \xi'_0&=A\xi_0-B\xi_1-C\xi_2-D\xi_3,\\
 \xi'_1&=B\xi_0+A\xi_1-D\xi_2+C\xi_3,\\
 \xi'_2&=C\xi_0+D\xi_1+A\xi_2-B\xi_3,\\
 \xi'_3&=D\xi_0-C\xi_1+B\xi_2+A\xi_3.
\end{aligned}
\]
由此可知，若考察一个平面 $(\xi)$，它通过点 $(x)$ 与其新位置 $(x')$ 的连线，则在滑移 (III) 以后，该平面的新位置 $(\xi')$ 仍通过连线 $(x)(x')$。事实上，由假设
\[
 \mathop{\mathrm{S}}x\xi=0,
 \qquad \mathop{\mathrm{S}}x'\xi=0
\]
利用 (III)、(III) 立即得到
\[
 \mathop{\mathrm{S}}x'\xi'=0,
 \qquad \mathop{\mathrm{S}}x\xi'=0.
\]
再注意，对应两平面的夹角 $p$ 满足
\[
 \cos p=\mathop{\mathrm{S}}\xi_i\xi'_i
\]
而由 (III) 得
\[
 \cos p=A,
\]
即
\[
 p=\frac\delta R.
\]
因此得到：\emph{在每个滑移中，连接两个对应点的直线在自身上滑动；这条连线上的每一点沿连线本身移动常长度 $\delta=R\arc\cos A$，而通过连线的每个平面都绕该直线转过常角 $p=\frac\delta R$。}

若要考察使固定点 $(1,0,0,0)$ 沿同一直线运动的单参数滑移族，只须令
\[
 A=\cos t,
 \qquad B=\beta\sen t,
 \qquad C=\gamma\sen t,
 \qquad D=\delta\sen t
\]
其中 $t$ 是变参数，而 $\beta,\gamma,\delta$ 是满足
\[
 \beta^2+\gamma^2+\delta^2=1.
\]
的三个常数。公式 (III) 或 (III*) 定义一个单参数连续滑移群；在空间的这种连续滑移中，各点的轨迹都是直线，形成一个称为\emph{Clifford 直线汇}的直线汇。把滑移分成 (III)、(III*) 两类，对应于平面在相应连续滑移中向右或向左旋转的方向；因此，这两类滑移也称为\emph{右旋滑移或左旋滑移}。
